\section{Rank}
For an $m\times n$ matrix, the rank $r$ is the number of pivots, the
dimension of its column space, and the dimension of its row space. The
rank--nullity equation is
\[
 \dim\Nul A=n-r.
\]
Notice the number of \emph{columns} on the right. Row operations preserve
the row space and null space, but may change the column space. Consequently
use pivot columns of the \emph{original} matrix for a column-space basis;
use nonzero rows of an echelon form for a row-space basis. A basis for a
null space comes from solving the homogeneous equations with free variables.

\practice{4.6.P1}{Rank and nullity from the given echelon form}
The given echelon matrix $B$ has two nonzero rows and two pivot positions.
Thus $\rank A=2$. Since $A$ has five columns, rank--nullity gives
$\dim\Nul A=5-2=3$.

\practice{4.6.P2}{Choose the correct rows and columns}
The pivot columns of $B$ are columns $1,2$. Therefore a column-space basis
uses the original columns $(2,1,-7,4)$ and $(-1,-2,8,-5)$ of $A$.
A row-space basis uses the two nonzero rows of $B$:
\[
 (1,-2,-4,3,-2),\qquad(0,3,9,-12,12).
\]
The former vectors have four entries and live in $\R^4$; the latter have
five entries and live in $\R^5$. This is a useful check on their roles.

\practice{4.6.P3}{The next step toward a null-space basis}
Continue reducing $B$ to reduced echelon form: divide its second row by
$3$, then add twice that new row to its first. The nonzero rows become
$(1,0,2,-5,6)$ and $(0,1,3,-4,4)$. From there take $x_3,x_4,x_5$ free,
express $x_1,x_2$ in terms of them, and separate their coefficient vectors.
It is the free variables, not the pivot columns, that directly produce this basis.

\practice{4.6.P4}{Transpose does not change rank}
The rank theorem gives $\rank A^T=\rank A=2$. Therefore every row echelon
form of $A^T$ has exactly two pivot columns.

\exercise{4.6.1}{Three bases from the supplied echelon form}
The supplied $B$ has 2 pivots, so $\rank A=2$ and $\dim\Nul A=4-2=2$, without row reducing $A$. Its pivot columns are 1, 2.
\step{Column and row bases use different source matrices.}
A matrix whose columns form a basis of $\Col A$ is
\[ C=\begin{pmatrix}1&-4\\-1&2\\5&-6\end{pmatrix}.\]
These are the pivot columns of $A$, not $B$. A matrix whose rows form a basis of $\operatorname{Row}A$ is
\[ R=\begin{pmatrix}1&0&-1&5\\0&-2&5&-6\end{pmatrix}.\]
\step{Solve the homogeneous equations for the free variables.}
Back substitution in $Bx=0$ gives a basis matrix
\[ N=\begin{pmatrix}2&-5\\5&-3\\2&0\\0&1\end{pmatrix},\qquad\Nul A=\{Nz:z\in\R^{2}\}.\]
The columns of $N$ are the coefficient vectors of the free parameters (some have been multiplied by nonzero scalars to avoid fractions). Each solves $Bx=0$, and hence $Ax=0$. They are independent because their free-coordinate entries isolate the parameters; their number equals the number of free variables. Thus they span the entire null space, not only some solutions.

\exercise{4.6.2}{Three bases from the supplied echelon form}
The supplied $B$ has 3 pivots, so $\rank A=3$ and $\dim\Nul A=5-3=2$, without row reducing $A$. Its pivot columns are 1, 3, 5.
\step{Column and row bases use different source matrices.}
A matrix whose columns form a basis of $\Col A$ is
\[ C=\begin{pmatrix}1&4&2\\2&6&-3\\3&3&-3\\3&0&0\end{pmatrix}.\]
These are the pivot columns of $A$, not $B$. A matrix whose rows form a basis of $\operatorname{Row}A$ is
\[ R=\begin{pmatrix}1&3&4&-1&2\\0&0&1&-1&1\\0&0&0&0&-5\end{pmatrix}.\]
\step{Solve the homogeneous equations for the free variables.}
Back substitution in $Bx=0$ gives a basis matrix
\[ N=\begin{pmatrix}-3&-3\\1&0\\0&1\\0&1\\0&0\end{pmatrix},\qquad\Nul A=\{Nz:z\in\R^{2}\}.\]
The columns of $N$ are the coefficient vectors of the free parameters (some have been multiplied by nonzero scalars to avoid fractions). Each solves $Bx=0$, and hence $Ax=0$. They are independent because their free-coordinate entries isolate the parameters; their number equals the number of free variables. Thus they span the entire null space, not only some solutions.

\exercise{4.6.3}{Three bases from the supplied echelon form}
The supplied $B$ has 3 pivots, so $\rank A=3$ and $\dim\Nul A=6-3=3$, without row reducing $A$. Its pivot columns are 1, 2, 5.
\step{Column and row bases use different source matrices.}
A matrix whose columns form a basis of $\Col A$ is
\[ C=\begin{pmatrix}2&6&3\\-2&-3&0\\4&9&3\\-2&3&3\end{pmatrix}.\]
These are the pivot columns of $A$, not $B$. A matrix whose rows form a basis of $\operatorname{Row}A$ is
\[ R=\begin{pmatrix}2&6&-6&6&3&6\\0&3&0&3&3&0\\0&0&0&0&3&0\end{pmatrix}.\]
\step{Solve the homogeneous equations for the free variables.}
Back substitution in $Bx=0$ gives a basis matrix
\[ N=\begin{pmatrix}3&0&-3\\0&-1&0\\1&0&0\\0&1&0\\0&0&0\\0&0&1\end{pmatrix},\qquad\Nul A=\{Nz:z\in\R^{3}\}.\]
The columns of $N$ are the coefficient vectors of the free parameters (some have been multiplied by nonzero scalars to avoid fractions). Each solves $Bx=0$, and hence $Ax=0$. They are independent because their free-coordinate entries isolate the parameters; their number equals the number of free variables. Thus they span the entire null space, not only some solutions.

\exercise{4.6.4}{Three bases from the supplied echelon form}
The supplied $B$ has 5 pivots, so $\rank A=5$ and $\dim\Nul A=6-5=1$, without row reducing $A$. Its pivot columns are 1, 2, 3, 5, 6.
\step{Column and row bases use different source matrices.}
A matrix whose columns form a basis of $\Col A$ is
\[ C=\begin{pmatrix}1&1&-2&1&-2\\1&2&-3&-2&-3\\1&-1&0&1&6\\1&-2&2&-3&0\\1&-2&1&2&-1\end{pmatrix}.\]
These are the pivot columns of $A$, not $B$. A matrix whose rows form a basis of $\operatorname{Row}A$ is
\[ R=\begin{pmatrix}1&1&-2&0&1&-2\\0&1&-1&0&-3&-1\\0&0&1&1&-13&-1\\0&0&0&0&1&-1\\0&0&0&0&0&1\end{pmatrix}.\]
\step{Solve the homogeneous equations for the free variables.}
Back substitution in $Bx=0$ gives a basis matrix
\[ N=\begin{pmatrix}-1\\-1\\-1\\1\\0\\0\end{pmatrix},\qquad\Nul A=\{Nz:z\in\R^{1}\}.\]
The columns of $N$ are the coefficient vectors of the free parameters (some have been multiplied by nonzero scalars to avoid fractions). Each solves $Bx=0$, and hence $Ax=0$. They are independent because their free-coordinate entries isolate the parameters; their number equals the number of free variables. Thus they span the entire null space, not only some solutions.

\exercise{4.6.5}{A rank-three matrix with seven columns}
Rank--nullity gives $\dim\Nul A=7-3=4$. The row-space dimension equals
rank, so $\dim\operatorname{Row}A=3$. Transpose preserves rank, so $\rank A^T=3$.

\exercise{4.6.6}{A rank-two matrix with five columns}
There are $5-2=3$ nonpivot columns, giving $\dim\Nul A=3$.
The rank theorem gives $\dim\operatorname{Row}A=2$ and $\rank A^T=2$.

\exercise{4.6.7}{Full row rank does not identify a null space with its coordinate space}
Four pivot columns in a $4\times7$ matrix mean rank $4$, so its columns span
all of $\R^4$: $\Col A=\R^4$. Rank--nullity gives nullity $7-4=3$.
But $\Nul A$ is a three-dimensional subspace of $\R^7$, not the set
$\R^3$ of triples. Equality $\Nul A=\R^3$ is therefore false; a choice of
basis gives an isomorphism with $\R^3$, not literal equality of the sets.

\exercise{4.6.8}{Rank four with six output coordinates}
The nullity is $8-4=4$. The column space is a four-dimensional subspace of
$\R^6$, because the columns have six entries. Thus it is not literally
$\R^4$; dimension does not change the number of entries in its vectors.

\exercise{4.6.9}{Recover rank from a known nullity}
The matrix has six columns and nullity three, so its rank is $6-3=3$.
Thus $\dim\Col A=3$. But the columns have four entries, so $\Col A$ is a
subspace of $\R^4$, not literally $\R^3$.

\exercise{4.6.10}{Seven columns and five free coordinates}
By rank--nullity, $\rank A=7-5=2$. Column-space dimension equals rank,
so the requested dimension is $2$. The eight rows specify the ambient
output space but do not replace the seven in the rank--nullity equation.

\exercise{4.6.11}{From nullity to row-space dimension}
Here $n=5$ and nullity is $3$, so rank is $5-3=2$.
The row-space dimension equals rank and is therefore $2$.

\exercise{4.6.12}{A second row-space dimension}
There are four columns and nullity $2$. Hence rank, and therefore row-space
dimension, equals $4-2=2$.

\exercise{4.6.13}{The largest possible rank}
There can be no more pivots than rows or columns, so an $m\times n$ matrix
has rank at most $\min(m,n)$. Both a $7\times5$ and a $5\times7$ matrix
therefore have largest possible rank $5$. This is attainable by placing a
$5\times5$ identity matrix in the first five rows or columns and filling
remaining entries with zeros.

\exercise{4.6.14}{The largest row-space dimension}
The row-space dimension is rank. For sizes $5\times4$ and $4\times5$,
the maximum rank is $4$ in each case. Identity blocks with an additional
zero row or zero column show that this bound is attained.

\exercise{4.6.15}{The smallest nullity for a wide matrix}
A $3\times7$ matrix has rank at most $3$, so its nullity is
$7-\rank A\ge7-3=4$. The matrix $[\,I_3\ 0\,]$ attains rank $3$ and
nullity $4$. Thus the smallest possible nullity is $4$.

\exercise{4.6.16}{The smallest nullity for a tall matrix}
A $7\times5$ matrix can have rank $5$, so its nullity can be $5-5=0$.
For example stack two zero rows below $I_5$. Nullity cannot be negative,
so zero is the minimum.

\exercise{4.6.17}{Five rank statements}
\begin{enumerate}[label=(\alph*)]
\item \textbf{True.} The columns of $A^T$ are exactly the transposed rows
of $A$; their spans are the same coordinate vectors.
\item \textbf{False.} The first original rows need not be independent.
For $A=\left(\begin{smallmatrix}1&0&0\\2&0&0\\0&1&0\\0&0&1\end{smallmatrix}\right)$
there are three pivots but the first three rows are dependent.
\item \textbf{True.} Both dimensions equal rank, including for rectangular matrices.
\item \textbf{False.} The sum equals the number of \emph{columns}, $n$,
not the number of rows $m$.
\item \textbf{True.} Finite-precision rounding can turn an exact zero
pivot into a small nonzero value or the reverse. Exact elementary operations
preserve rank; its apparent computed value may change, as the Numerical Note explains.
\end{enumerate}

\exercise{4.6.18}{Five statements about row operations}
\begin{enumerate}[label=(\alph*)]
\item \textbf{False.} Use pivot columns of $A$ for $\Col A$, not generally
those of its echelon form. For example swapping the rows of
$\left(\begin{smallmatrix}1&0\\0&0\end{smallmatrix}\right)$ changes its column span.
\item \textbf{False.} The row span is preserved, but the particular
relations among the rows can change. Rows $(1,0),(1,0)$ satisfy $r_2-r_1=0$;
a row replacement gives rows $(1,0),(0,0)$, for which that relation no longer holds.
\item \textbf{True.} Nullity equals the number of free variables,
namely the number of nonpivot columns.
\item \textbf{True.} Transposing the rows of $A^T$ gives the columns of $A$.
\item \textbf{True.} Every new row is a combination of old rows, and
invertibility of elementary operations gives the reverse inclusion of row spans.
\end{enumerate}

\exercise{4.6.19}{A one-dimensional kernel of a five-row system}
All homogeneous solutions are multiples of one nonzero solution, so nullity
is $1$. With six columns, rank is $6-1=5$, equal to the number of rows.
Thus there is a pivot in every row, and $Ax=b$ is consistent for every
right-hand side $b\in\R^5$. The answer is yes.

\exercise{4.6.20}{Two free variables in an eight-variable system}
A consistent system with eight variables and two free variables has six
pivots. There are also six equations, so every row has a pivot. Therefore
$Ax=b$ is consistent for every $b\in\R^6$; changing constants alone cannot
make it inconsistent. The answer is no.

\exercise{4.6.21}{Onto with just one free direction}
Consistency for every right-hand side in a nine-row system means rank $9$.
With ten variables the nullity is $10-9=1$. Thus all homogeneous solutions
are multiples of one vector. It is impossible to find two nonzero
homogeneous solutions that are not multiples of each other.

\exercise{4.6.22}{Too many variables for a one-dimensional kernel}
For ten equations in twelve variables, rank is at most ten. Nullity is
therefore at least $12-10=2$. It cannot be $1$, so not all homogeneous
solutions can be multiples of one fixed nonzero vector.

\exercise{4.6.23}{Replace redundant equations}
The two given nonproportional solutions are independent and span the kernel,
so nullity is $2$. There are eight variables, giving rank $8-2=6$.
Row reduction therefore leaves six nonzero rows and six redundant zero rows.
Keeping just the six nonzero homogeneous equations describes exactly the same
solutions. Fewer than six cannot suffice: a matrix with fewer rows would
have rank below six and nullity above two.

\exercise{4.6.24}{A unique solution for one right-hand side versus all}
A $7\times6$ matrix may have rank $6$, giving zero kernel. Then every
consistent system has a unique solution; for example
$A=\left(\begin{smallmatrix}I_6\\0\end{smallmatrix}\right)$ and $b=0$ do.
But rank is at most $6<7$, so the column space cannot be all of $\R^7$.
Thus a unique solution for some right-hand side is possible, while a unique
solution for every right-hand side is impossible.

\exercise{4.6.25}{One missing pivot row leaves some constants unattainable}
Twelve variables and three free variables mean rank $9$. The system has ten
rows, so it lacks a pivot in one row. Its column space is a proper subspace
of $\R^{10}$, and some right-hand sides lie outside it. Thus the scientist
cannot guarantee a solution after changing the constants; indeed there
necessarily exist choices of constants that make the system inconsistent.

\exercise{4.6.26}{Full rank for a tall matrix}
When $m>n$, the largest possible rank is $n$. The columns are independent
if and only if all $n$ columns are pivot columns, which is equivalent to
rank $n$. This is precisely full rank for the given shape of matrix.

\exercise{4.6.27}{Locate the four fundamental spaces}
Rows of $A$ and inputs to $A$ have $n$ entries; columns of $A$ and inputs
to $A^T$ have $m$ entries. Consequently
\[
 \begin{array}{c|c}
 \R^n&\operatorname{Row}A=\Col A^T,\quad\Nul A\\
 \R^m&\Col A=\operatorname{Row}A^T,\quad\Nul A^T
 \end{array}
\]
The six names reduce to four fundamental spaces after the two displayed
identifications. For special matrices some of these four may coincide as
sets (for example several can be the zero space); ``four'' refers to the
four roles, not a guarantee of four different sets in every example.

\exercise{4.6.28}{Two rank--nullity identities}
\step{(a) Apply rank--nullity to $A$.}
Since $\dim\operatorname{Row}A=\rank A$ and
$\dim\Nul A=n-\rank A$, their sum is $n$.
\step{(b) Apply rank--nullity to $A^T$.}
This transpose has $m$ columns and rank equal to that of $A$.
Thus $\dim\Nul A^T=m-\rank A$; adding $\dim\Col A=\rank A$ gives $m$.

\exercise{4.6.29}{Onto is equivalent to a trivial transpose kernel}
The equation $Ax=b$ is solvable for every $b\in\R^m$ if and only if
$\dim\Col A=m$. By Exercise 4.6.28(b), this is equivalent to
$\dim\Nul A^T=0$, which is equivalent to $\Nul A^T=\{0\}$.
That last condition says exactly that $A^Tx=0$ has only the trivial solution.

\exercise{4.6.30}{The augmented-rank consistency criterion}
Adding $b$ as a new column increases rank precisely when $b$ is outside
the span of the old columns. Therefore $Ax=b$ is consistent if and only if
\[
 \rank[\,A\ b\,]=\rank A.
\]
Equivalently, an inconsistent row in the reduced augmented matrix gives an
extra pivot in its last column. If ranks agree, that extra pivot is absent.

\exercise{4.6.31}{An outer product has at most one independent column}
For $u=(2,-3,5)$ and $v=(a,b,c)$,
\[
 uv^T=\begin{pmatrix}2a&2b&2c\\-3a&-3b&-3c\\5a&5b&5c\end{pmatrix}.
\]
Its columns are $au,bu,cu$, so its column space is contained in
$\Span\{u\}$, a one-dimensional space. Hence rank is at most $1$.
It is zero when $v=0$, and otherwise is exactly $1$, since $u\ne0$.

\exercise{4.6.32}{Factor a rank-one matrix}
With $u=(1,2)$ the matrix $uv^T$ has first row $v^T$ and second row $2v^T$.
Thus choose $v=(1,-3,4)$, giving
\[
 \begin{pmatrix}1\\2\end{pmatrix}[\,1\ -3\ 4\,]
 =\begin{pmatrix}1&-3&4\\2&-6&8\end{pmatrix}.
\]

\exercise{4.6.33}{Every nonzero column generates a rank-one column space}
If rank is $1$ and the first column $u$ is nonzero, it is a basis for the
column space. Thus the three columns have form $u,\alpha u,\beta u$.
Taking $v=(1,\alpha,\beta)$ gives $A=uv^T$.
If the first column is zero, choose any other nonzero column as $u$ (one
exists because rank is $1$). Express each column as $v_j u$, including
coefficient zero for a zero column, and set $v=(v_1,v_2,v_3)$.
The same outer-product formula then holds.

\exercise{4.6.34}{A rank-$r$ matrix is a sum of $r$ rank-one matrices}
Since $U$ is obtained by elementary row operations, $U=FA$ for an invertible
product $F$ of elementary matrices. Set $E=F^{-1}$ to get $A=EU$.
Write $E=[\,e_1\ \cdots\ e_m\,]$ and let $u_i^T$ be row $i$ of $U$.
Matrix multiplication by columns and rows gives
\[
 A=EU=\sum_{i=1}^m e_i u_i^T=\sum_{i=1}^r e_i u_i^T.
\]
The last equality uses the zero rows below the first $r$ rows of the echelon
matrix. Each of those first $r$ rows is nonzero, and each column $e_i$ of
an invertible matrix is nonzero. Thus every displayed outer product has
rank exactly $1$, proving the assertion.

\exercise{4.6.35}{Four fundamental bases and two invertible assembled matrices}
\[ A=\begin{pmatrix}7&-9&-4&5&3&-3&-7\\-4&6&7&-2&-6&-5&5\\5&-7&-6&5&-6&2&8\\-3&5&8&-1&-7&-4&8\\6&-8&-5&4&4&9&3\end{pmatrix}.\]
\step{(a) Supply and verify the three requested bases.}
The pivot columns are 1, 2, 4, 6, so rank is $4$. The following matrices give respectively a column-space basis in their columns, a null-space basis in their columns, and a row-space basis in their rows:
\[ C=\begin{pmatrix}7&-9&5&-3\\-4&6&-2&-5\\5&-7&5&2\\-3&5&-1&-4\\6&-8&4&9\end{pmatrix}.\]
\[ N=\begin{pmatrix}-13&-10&3\\-11&-1&-2\\2&0&0\\0&11&-7\\0&2&0\\0&0&-1\\0&0&1\end{pmatrix}.\]
\[ R=\begin{pmatrix}1&0&\frac{13}{2}&0&5&0&-3\\0&1&\frac{11}{2}&0&\frac{1}{2}&0&2\\0&0&0&1&- \frac{11}{2}&0&7\\0&0&0&0&0&1&1\end{pmatrix}.\]
The rows of $R$ are the nonzero rows of the reduced echelon form. The $3$ columns of $N$ result from its free-variable equations; exact multiplication gives $AN=0$, and their free coordinates show independence. Thus these are bases, not only spanning lists.
\step{(b) Supply the transpose-kernel basis and check invertibility.}
\[ M=\begin{pmatrix}2\\41\\0\\-28\\11\end{pmatrix}.\]
Here $A^TM=0$, and its $1$ columns are independent. Form $S=[\,R^T\ N\,]$ and $T=[\,C\ M\,]$ using the displayed blocks. There are $4+3=7$ columns in $S\in M_{7\times 7}$ and $4+1=5$ columns in $T\in M_{5\times 5}$. Exact determinants are
\[ \det S=- \frac{88369}{2},\qquad\det T=10360.\]
Both are nonzero, verifying that both square matrices are invertible. Conceptually, a vector in both $\operatorname{Row}A$ and $\Nul A$ has form $x=A^Ty$ and satisfies $Ax=0$, so $x^Tx=y^TAx=0$. Since $x^Tx$ is a sum of squares, $x=0$. The same argument with $A^T$ treats $\Col A$ and $\Nul A^T$; their dimensions add to the full ambient dimensions. This explains why assembling their bases produces invertible matrices.

\exercise{4.6.36}{Four fundamental bases and two invertible assembled matrices}
The problem permits a generated example. Here is one reproducible sample: with Python \texttt{random.Random(46)}, draw the entries of a $6\times2$ matrix $J$ and then a $2\times7$ matrix $K$ successively from \texttt{randint(-1,1)}, and set $A=JK$. This guarantees rank at most $2$, which in particular is at most $4$. The resulting matrix is displayed below; its first two columns are independent, so its actual rank is $2$.
\[ A=\begin{pmatrix}1&1&0&-1&1&1&1\\2&1&1&-1&2&0&1\\-2&-1&-1&1&-2&0&-1\\0&-1&1&1&0&-2&-1\\2&1&1&-1&2&0&1\\-2&-1&-1&1&-2&0&-1\end{pmatrix}.\]
\step{(a) Supply and verify the three requested bases.}
The pivot columns are 1, 2, so rank is $2$. The following matrices give respectively a column-space basis in their columns, a null-space basis in their columns, and a row-space basis in their rows:
\[ C=\begin{pmatrix}1&1\\2&1\\-2&-1\\0&-1\\2&1\\-2&-1\end{pmatrix}.\]
\[ N=\begin{pmatrix}-1&0&-1&1&0\\1&1&0&-2&-1\\1&0&0&0&0\\0&1&0&0&0\\0&0&1&0&0\\0&0&0&1&0\\0&0&0&0&1\end{pmatrix}.\]
\[ R=\begin{pmatrix}1&0&1&0&1&-1&0\\0&1&-1&-1&0&2&1\end{pmatrix}.\]
The rows of $R$ are the nonzero rows of the reduced echelon form. The $5$ columns of $N$ result from its free-variable equations; exact multiplication gives $AN=0$, and their free coordinates show independence. Thus these are bases, not only spanning lists.
\step{(b) Supply the transpose-kernel basis and check invertibility.}
\[ M=\begin{pmatrix}0&2&0&0\\1&-1&-1&1\\1&0&0&0\\0&1&0&0\\0&0&1&0\\0&0&0&1\end{pmatrix}.\]
Here $A^TM=0$, and its $4$ columns are independent. Form $S=[\,R^T\ N\,]$ and $T=[\,C\ M\,]$ using the displayed blocks. There are $2+5=7$ columns in $S\in M_{7\times 7}$ and $2+4=6$ columns in $T\in M_{6\times 6}$. Exact determinants are
\[ \det S=23,\qquad\det T=-21.\]
Both are nonzero, verifying that both square matrices are invertible. Conceptually, a vector in both $\operatorname{Row}A$ and $\Nul A$ has form $x=A^Ty$ and satisfies $Ax=0$, so $x^Tx=y^TAx=0$. Since $x^Tx$ is a sum of squares, $x=0$. The same argument with $A^T$ treats $\Col A$ and $\Nul A^T$; their dimensions add to the full ambient dimensions. This explains why assembling their bases produces invertible matrices.

\exercise{4.6.37}{Reconstruct the original matrix from pivot columns}
Use the $C$ and $R$ displayed in Exercise 4.6.35. Those are precisely the
pivot columns of the original $A$ and the nonzero rows of its reduced
echelon form. Multiplication gives
\[
 CR=\begin{pmatrix}
 7&-9&-4&5&3&-3&-7\\-4&6&7&-2&-6&-5&5\\
 5&-7&-6&5&-6&2&8\\-3&5&8&-1&-7&-4&8\\6&-8&-5&4&4&9&3
 \end{pmatrix}=A.
\]
For example, the third column of $R$ is $(13/2,11/2,0,0)$, so the third
column of $CR$ is $(13/2)a_1+(11/2)a_2=a_3$.
The reduced matrix records the coefficients of every original column in
terms of the pivot columns. Thus this product restores all columns exactly.

\exercise{4.6.38}{Three generated rank factorizations and a general proof}
For a reproducible family, use \texttt{random.Random(4600+r)} for each
$r=5,4,3$. Form $C$ by stacking $I_r$ over a $(5-r)\times r$ array of
successive \texttt{randint(-2,2)} draws, and form $R$ by placing an
$r\times(7-r)$ array of subsequent draws after $I_r$. Set $A=CR$.
The first $r$ columns of $A$ are $C$; its reduced echelon form has nonzero
rows $R$. Thus these are actual examples of the exercise's construction,
with exact rank $r$, rather than uncontrolled random ranks.
\step{The rank-5 sample.}
\[ C=\begin{pmatrix}1&0&0&0&0\\0&1&0&0&0\\0&0&1&0&0\\0&0&0&1&0\\0&0&0&0&1\end{pmatrix},\qquad R=\begin{pmatrix}1&0&0&0&0&2&-1\\0&1&0&0&0&-1&-2\\0&0&1&0&0&0&2\\0&0&0&1&0&0&0\\0&0&0&0&1&-1&0\end{pmatrix}.\]
\[ CR=A=\begin{pmatrix}1&0&0&0&0&2&-1\\0&1&0&0&0&-1&-2\\0&0&1&0&0&0&2\\0&0&0&1&0&0&0\\0&0&0&0&1&-1&0\end{pmatrix}.\]
\step{The rank-4 sample.}
\[ C=\begin{pmatrix}1&0&0&0\\0&1&0&0\\0&0&1&0\\0&0&0&1\\-2&0&2&1\end{pmatrix},\qquad R=\begin{pmatrix}1&0&0&0&-1&-1&-1\\0&1&0&0&-2&-1&-1\\0&0&1&0&-2&0&-1\\0&0&0&1&-1&0&2\end{pmatrix}.\]
\[ CR=A=\begin{pmatrix}1&0&0&0&-1&-1&-1\\0&1&0&0&-2&-1&-1\\0&0&1&0&-2&0&-1\\0&0&0&1&-1&0&2\\-2&0&2&1&-3&2&2\end{pmatrix}.\]
\step{The rank-3 sample.}
\[ C=\begin{pmatrix}1&0&0\\0&1&0\\0&0&1\\2&2&-2\\1&2&2\end{pmatrix},\qquad R=\begin{pmatrix}1&0&0&1&0&-2&0\\0&1&0&2&-2&2&-1\\0&0&1&2&-1&-2&1\end{pmatrix}.\]
\[ CR=A=\begin{pmatrix}1&0&0&1&0&-2&0\\0&1&0&2&-2&2&-1\\0&0&1&2&-1&-2&1\\2&2&-2&2&-2&4&-4\\1&2&2&9&-6&-2&0\end{pmatrix}.\]

\step{Prove the observed equality for every matrix.}
Let $C$ consist of the original pivot columns of a rank-$r$ matrix $A$,
and $R$ of the nonzero rows of its reduced echelon form $U$.
In $U$, the pivot columns are the first $r$ coordinate vectors above the
zero rows. Therefore the entries of column $j$ of $R$ express column $j$
of $U$ as a combination of its pivot columns. Elementary row operations
preserve all column dependence relations: if $U=EA$ with $E$ invertible,
then $U\alpha=0$ if and only if $A\alpha=0$. The same coefficients thus
express column $j$ of $A$ using columns of $C$. This means column $j$ of
$CR$ is column $j$ of $A$, for every $j$; hence $CR=A$.
